\section{Certificate conventions}
\label{app:certificates}

\subsection{Exact polynomials}

The transport representation of a rational is the pair
\code{[numerator, denominator]}, with positive denominator and relatively prime
entries. A polynomial is a list of pairs \code{[exponentVector, rational]}.
The zero polynomial is the empty list. Exponents are nonnegative integers;
zero coefficients and repeated monomials are rejected. The variable arity is
supplied by the externally owned problem, not inferred from a certificate.

For instance, in variable order $(x,y)$, $\tfrac12x^2-3y$ is
\begin{lstlisting}
[
  [[2, 0], [1, 2]],
  [[0, 1], [-3, 1]]
]
\end{lstlisting}
The checker need not require a particular monomial ordering for acceptance.
The searcher uses a deterministic order for row reduction and graded
lexicographic order for its Gr\"obner routine. These are algorithmic choices,
not mathematical assumptions about the polynomial identities.

\subsection{Action records}

For \code{pullback-space-v1}, \code{basis} is a list of $r$ polynomials,
\code{actions[a][i][j]} is a rational scalar, and
\code{target\_coefficients[i]} is a rational scalar. For
\code{pullback-ideal-v1}, \code{generators} replaces \code{basis}, and the
entries of \code{actions} and \code{target\_coefficients} are polynomials.
The former is the constant-multiplier special case of the latter, but keeping
separate tags makes malformed scalar/polynomial payloads easier to reject.

The problem owns initialization, transitions, and the target. The certificate
cannot change the target by returning a different convenient equation. Empty
or redundant invariant families may be valid, but they prove a nonzero target
only when the checked target combination actually reconstructs it.

\subsection{Ranking records}

An affine expression is encoded constant-first:
\[
 [b,w_1,\ldots,w_n]\quad\longleftrightarrow\quad b+\sum_iw_ix_i.
\]
Every \code{layers[k]} maps every node tag to its affine rank at that layer.
Every edge has one witness record with its pivot, a decrease combination, and
one nonnegativity combination per updated coordinate. Multipliers are ordered
against explicit edge guards, then $x_1,\ldots,x_n$, then $1$.

The two-phase example uses
\begin{lstlisting}
"layers": [ { "P": [1, 2], "Q": [0, 2] } ]
\end{lstlisting}
The checker verifies all edges in their original order. Dropping a hard edge
from the certificate is rejected by the coverage check; adding another easy
edge's witness in its place fails the exact coefficient identity.

\subsection{Orbit records}

An orbit counterexample contains rational initialization parameters and a list
of transition indices in \emph{execution order}. It is checked by direct exact
evaluation of the problem's maps. For the noncommuting fixture,
\[
 F_0(x,y)=(y,0),\qquad F_1(x,y)=(0,1),\qquad (x,y)_0=(0,0),
\]
the target is $x=0$. Word $[1,0]$ reaches $(1,0)$ and violates it; word $[0,1]$
ends at $(0,1)$ and does not. This is why reversing the word is a meaningful
corruption rather than another valid counterexample.

\section{Why the finite synthesis counts have the reported shape}

With $n$ distractor element binders, the four hole grammars have sizes
$n+1,n+2,n+2,n+8$. The first locally valid base is last, the valid negative
branch is last, the valid zero branch is second-to-last, and the valid positive
continuation application is last. In the fixed tuple order, the successful
complete term therefore has ordinal
\[
 (n+1)(n+2)^2(n+8)-(n+8).
\]
For $n=12$ this is $50\,940$. Local filtering tests every component once, so its
component count is
\[
 (n+1)+2(n+2)+(n+8)=4n+13,
\]
and each hole retains one candidate, yielding one assembled term.

The contrast is quartic versus linear in this deliberately controlled family.
It is not a complexity theorem for unrestricted synthesis. Its assumptions are
visible: fixed sketch, independent exact local contracts, finite grammars,
distinct distractors, and the stated order. Correlated holes need a join, and
local contract discovery itself can be expensive.

\section{Reproducing the package}

From the archive root, use:
\begin{lstlisting}
python -S -m unittest discover -s prototype -p "test_*.py" -v
python -S prototype/run_experiments.py --out reproduced-results
python -S prototype/replay.py reproduced-results/certificates.json
\end{lstlisting}
The optional independent CAS comparison is:
\begin{lstlisting}
python prototype/validate_sympy.py --out reproduced-results/sympy-validation.json
\end{lstlisting}
It requires SymPy; the recorded comparison used version 1.14.0. The other commands
do not. The archived \code{results/} directory contains the original evidence,
while reruns default to \code{reproduced-results/} to avoid overwriting it.

To rebuild this article:
\begin{lstlisting}
cd article
pdflatex -interaction=nonstopmode -halt-on-error main.tex
pdflatex -interaction=nonstopmode -halt-on-error main.tex
pdflatex -interaction=nonstopmode -halt-on-error main.tex
\end{lstlisting}
The distributed PDF is named \code{forge-extensions.pdf}. A fresh build produces
\code{main.pdf}; the build script copies it to the distributed name. The required
LaTeX packages are listed by the preamble. No font files are distributed.

\subsection{A first Lean verification session}

In a project with the repository's pinned Lean toolchain, the first experiment
should be to compile \code{lean/Specimens.lean}, inspect the printed axiom sets,
and retain the compiler output as a new receipt. This archive supplies no such
receipt. Next, lower one small invariant certificate to a theorem using the
pinned Mathlib environment and prove the source reification lemma for that exact
fixture. Only then count it as a Lean acceptance.

The proof specimens cover fold correctness for an explicit list encoding,
invariant propagation along a reachability derivation, and closing recursive
calls by a well-founded induction hypothesis. They deliberately do not define
a pretend \forge{} tactic that succeeds by calling an unimplemented backend.
